\documentclass[UTF8,a4paper,11pt,fontset=fandol]{ctexart}
\usepackage[margin=23mm,headheight=15pt]{geometry}
\usepackage{amsmath,amssymb,amsthm,mathtools,bm}
\usepackage{booktabs,longtable,tabularx,array,enumitem}
\usepackage{xurl}
\usepackage[hidelinks,bookmarksnumbered]{hyperref}
\usepackage{fancyhdr}
\pagestyle{fancy}\fancyhf{}\fancyhead[L]{Yau26\quad 研究续接基线·覆盖补全v2}\fancyhead[R]{2026-09-16}\fancyfoot[C]{\thepage}
\setlength{\parskip}{4pt}\linespread{1.13}
\setlength{\emergencystretch}{3em}
\setlist{itemsep=2pt,topsep=4pt,parsep=0pt,leftmargin=2em}
\setcounter{tocdepth}{1}
\numberwithin{equation}{section}
\newtheorem{theorem}{定理}[section]
\newtheorem{proposition}[theorem]{命题}
\newtheorem{lemma}[theorem]{引理}
\newtheorem{target}[theorem]{研究目标}
\theoremstyle{definition}
\newtheorem{definition}[theorem]{定义}
\newtheorem{remark}[theorem]{说明}
\newcommand{\Ric}{\operatorname{Ric}}
\newcommand{\Rm}{\operatorname{Rm}}
\newcommand{\Vol}{\operatorname{Vol}}
\newcommand{\diam}{\operatorname{diam}}
\newcommand{\tr}{\operatorname{tr}}
\newcommand{\spec}{\operatorname{Spec}}
\newcommand{\Met}{\operatorname{Met}}
\newcommand{\Id}{\operatorname{Id}}
\newcommand{\op}{\mathrm{op}}
\newcommand{\CP}{\mathbb{CP}}
\newcommand{\R}{\mathbb R}
\newcommand{\Z}{\mathbb Z}
\newcommand{\dd}{\,d\mu_g}
\newcommand{\avgint}{\mathop{\rlap{\,\raisebox{0.1ex}{\tiny$-$}}\int}\nolimits}
\newcommand{\fint}{\mathop{\rlap{\,\raisebox{0.1ex}{\tiny$-$}}\int}\nolimits}
\newcommand{\norm}[1]{\left\lVert#1\right\rVert}
\newcommand{\source}[1]{\par\noindent\textit{来源与状态：}#1\par}
\newcommand{\asset}[2]{\par\noindent\textbf{#1}\quad\path{#2}\par}
\newcolumntype{L}[1]{>{\raggedright\arraybackslash}p{#1}}
\newcommand{\status}[1]{\textbf{【#1】}}

\title{主线三：近Einstein构造与全局拓扑的闭合\\[0.5em]{\large S0910A03｜图流形、曲面丛、全局存在性与Yau26的决定性接口}}
\author{Qiushi Engine\quad 科学续接稿}
\date{2026年9月16日\quad 覆盖补全版 v2}
\begin{document}
\maketitle
\begin{abstract}
本主线不是第三项已经完成的重大突破，而是连接两条已有成果主线与原题的全局研究线。它同时保留两类互补对象：已知无Einstein但尚未构造任意紧Ricci夹逼的extended graph流形；记录已有近负Einstein输入但尚未排除全部Einstein度量的曲面丛与其他固定解消。第100步的相对Chern--Gauss--Bonnet预算、完整矩阵连接器、周边格点与长颈参数联动是直接起点。本稿将其整理为明确的全局定理目标，并由文献中的最小Ricci体积重新推导图流形近夹逼序列必须满足的体积和Weyl能量饱和条件。
\end{abstract}
\begingroup
\small\linespread{1.0}\selectfont\setlength{\parskip}{0pt}
\tableofcontents
\endgroup
\clearpage

\section*{覆盖复核说明：本稿的边界与原始研究的优先性}
\addcontentsline{toc}{section}{覆盖复核说明}
本版为原三份续接稿的覆盖补全版（v2）。它\emph{不声称把全部原始证明和300步记录压缩成了无损替代物}。上一版覆盖了主要目标，但遗漏或过度压缩了部分最终成果、技术依赖和候选分支；本版补回具体条目，并配套保存可访问原件与逐步索引。正文中“完整”“已核查”等表述只能按紧邻的来源和范围理解，不能解释为对原工作区全部研究的认证。

必须区分三类内容：原始归档或轨迹记载的研究；用户转述的Codex原稿审读；上一版编辑者的代数补充、外部参考与路线建议。后两类不得替换前一类。原稿中尚未取得的系数、PDE细节、程序和历史题页，仍以\emph{原件待取得}记录，不能根据摘要补成已证内容。上一版新定义的密度或新推论属于可选校验材料，不是原研究新增了一项已经认证的成果。

A01/A02/A03的三份分工是编辑组织建议，不是历史上三个会话原有任务的逐字复制。先读取各自Step100的实际交接和当前工作区，再决定是否采用建议；任何原有路线不因未被安排为首要方向就被宣告失败或删除。总体目标仍是原始Yau26，不预定真或假，也不预定所有有效突破必然来自当前两条成果主线。

配套\path{审计/300步原始任务索引.json}逐条保存三份实际上传轨迹的300个步骤的\path{tx_task}、\path{rx_task}、状态及行号；\path{审计/研究条目覆盖表.tsv}列出本次人工对照的研究单元。\path{原始材料/}保留三份完整meta-trace及两份输入ZIP，字节校验见\path{审计/原件校验.tsv}。这些索引证明可追溯与原件保留，\emph{不证明每条历史主张正确，也不等于数学语义覆盖率100\%}。

旧ZIP引用的\path{meta_trace_report (93)(1).md}、\path{(94)(1).md}、\path{(95)(1).md}，与本次可见的\path{(100).md}及两份时间戳文件，尚无逐字同一性证明。三份现有轨迹按其内部工作区标记对应A01/A02/A03；不能只因都为100步就替换旧文件身份。归档另外引用的\path{Yau26_Ricci_pinching_research.md}、\path{Yau26_synthesis.tex/pdf}、\path{Yau26_uniform_gap.tex}和服务器工作区，不在当前可打包原件中。

\section{原题、量词与成功标准}
本会话组的工作命题保持为：对每个维数 $n\ge2$，是否存在只依赖 $n$ 的实常数 $c_n<C_n$，使任意闭、连通、光滑实流形 $M^n$ 只要承载一个度量 $g$ 满足
\begin{equation}\label{eq:Q}
 c_ng\le\Ric_g\le C_ng,
\end{equation}
就承载某个光滑 Einstein 度量 $h$，即 $\Ric_h=\lambda h$，$\lambda\in\R$？不要求 $h$ 接近 $g$，不要求 $h$ 与 $g$ 的 Ricci 符号相同；结论针对同一\emph{光滑流形}，不是仅同胚类型。

历史归属为 Yau 1982 年 \emph{Problem Section} 的第26题\cite{Yau1982}。当前归档没有附上原始题页的完整可核对影像；上式是三个会话反复确认的工作表述，而不是本文件宣称已经逐字认证的历史引文。续接时应在当前工作区核对原始任务书与题页，若发现实质差异，单列差异并调整对应推论，不把相邻的稳定性问题替换成原题。

定义
\begin{equation}\label{eq:defect}
 \delta_\sigma(g)=\norm{g^{-1}\Ric_g-\sigma\Id}_{L^\infty,\op},\qquad
 e_\sigma(M)=\inf_{g\in\Met(M)}\delta_\sigma(g),\qquad \sigma\in\{+1,-1\}.
\end{equation}
$P_n^\sigma$ 表示存在仅依赖 $n$ 的 $\varepsilon_n>0$，使 $\delta_\sigma(g)<\varepsilon_n$ 推出 $M$ 承载某个 Einstein 度量。工作命题满足
\begin{equation}\label{eq:branches}
 Q_n\quad\Longleftrightarrow\quad P_n^+\ \lor\ P_n^-.
\end{equation}
证明只用缩放：$\widehat g=ag$ 时，作为 $(0,2)$ 张量的 Ricci 不变，Ricci 端同态变为 $a^{-1}g^{-1}\Ric_g$。任一非退化区间都有一个非零内点 $\mu$ 及同号小邻域；将 $\Ric\approx\operatorname{sign}(\mu)g$ 的度量按 $1/|\mu|$ 缩放即可进入该邻域。反向直接取以 $\pm1$ 为中心的小区间。

因此，一个固定维数的\emph{证明}只需要完成其中一个符号分支；\emph{完整反驳}要在同一维数击败两个符号分支，可以使用不同流形。每个流形都必须排除任意符号的全部 Einstein 度量。构造 $e_\sigma(M)=0$ 本身是一项独立结论，但它只有与该 $M$ 的全部 Einstein 不存在性相结合，才构成分支反例。

二维由一致化、三维正分支由 Hamilton 的正 Ricci 结果提供经典基线；不得把这些作为本项目的新成果。四维的专门工具不会自动推广到 $n\ge5$。本文件允许严谨的强中间结果独立成文，但不把它们改称原题的完整解决。

\subsection{统一范数，避免四倍因子再次混入}
下列 $W^\pm$ 与签名公式在定向四维范畴内使用；非定向情形须另行处理定向双覆盖及等变下降，不能由覆盖的Einstein存在直接推出原流形结论。四维设 $s=\tr_g\Ric$、$E=\Ric-\frac{s}{4}g$。本文件的 $W^\pm$ 使用在 $\Lambda^2_\pm$ 上的曲率\emph{算子}范数；$E$ 和 Ricci 使用张量范数。于是
\begin{align}
8\pi^2\chi(M)&=\int_M\left(|W^+|^2+|W^-|^2+\frac{s^2}{24}-\frac{|E|^2}{2}\right)\dd,\label{eq:CGB}\\
12\pi^2\tau(M)&=\int_M(|W^+|^2-|W^-|^2)\dd,\label{eq:signature}\\
|\Rm_{\rm full}|^2&=4(|W^+|^2+|W^-|^2)+2|E|^2+\frac{s^2}{6}.
\end{align}
若写 $\mathcal R$ 表示完整曲率算子，则 $|\Rm_{\rm full}|^2=4|\mathcal R|^2$。精确 Einstein 时分别有 $\int|\mathcal R|^2=8\pi^2\chi$ 和 $\int|\Rm_{\rm full}|^2=32\pi^2\chi$。Eguchi--Hanson 的 $12\pi^2$ 能量采用前一种规范；后一种为 $48\pi^2$。

$\delta_\sigma\le\delta$ 给出
\begin{equation}
 |s-4\sigma|\le4\delta,\qquad |E|^2\le4\delta^2,\qquad
 |\Ric-\sigma g|^2\le4\delta^2.
\end{equation}
全文不把小 $L^p$ 积分、小加权范数、局部精确解，替换成原题要求的全域无权 Ricci 端同态上确界。

\section{材料边界与使用方式}
本文件是供原会话从 Step101 继续研究的科学续接稿，不是以整理文字替代原始证明的投稿稿。依据包括：2026-09-15进展包、2026-09-16逐步归档、S0910A01/A02/A03各100步的完整 meta-trace，以及用户提供的 Codex 对原稿的逐段审读报告。外部文献核对日期为2026-09-16。两个ZIP含摘要、公式、分级和原稿路径，但不含所有被引用的原始 TeX、证明笔记和程序字节。

下列标签只用于区分证据，不对数学内容按平台打分：\status{文献}表示本轮核对到的外部定理或明确引用；\status{原稿审读}表示原稿被 Codex 实质审读、未发现缺口的报告；\status{轨迹}表示会话记载的推导或内部核验；\status{本稿重算}表示下文完整展示的代数或逻辑推导；\status{待证}表示仍需新证明的目标。已撤回结论只作为边界保存，不进入定理依赖。

精确 Einstein 的统一\emph{拓扑}间隙已有原稿审读支持，不能再因一般 near-Einstein 平滑接口未闭合而降回“只有摘要的想法”。一般 near-Einstein 的统一\emph{Ricci夹逼}间隙仍未完成，不能因前一项成立就自动升级。固定 $A_1$ 近 Einstein 构造与全局 Einstein 不存在性也是两个命题。

第一次接续应读取本会话当前任务书、工作区最近修正及下文列出的原稿，记录版本；\emph{不重启已经做完的300步，也不把审读报告当成原稿全文}。三份文件彼此独立可编译；没有要求安装新应用或改变系统基础架构。缺失来源的恢复应服务于一个具体定理，不应演变为无期限重复检索。

\section{本主线的任务：让两个性质作用于同一光滑对象}
对一个明确指定的 $M$，反例所需的两部分是
\[
\mathsf P_\sigma(M):\ e_\sigma(M)=0,
\qquad
\mathsf N(M):\ M\text{ 不承载任何Einstein度量}.
\]
本会话的科学目标是严格建立某个 $\mathsf P_\sigma(M)\land\mathsf N(M)$，或者证明关键候选的 $e_\sigma$ 有正下界，或者反向构造该候选上的Einstein度量。三种结论都推进原题，但逻辑作用不同。

\begin{longtable}{L{.22\textwidth}L{.34\textwidth}L{.35\textwidth}}
\toprule 对象 & 当前可以继承的内容 & 真正尚未完成的内容\\\midrule\endhead
指定一般型吹起族 & A01持有精确拓扑间隙及必要退化的审读支持 & 近似方程的无额外几何假设间隙，或同一无Einstein拓扑上的输入构造\\
固定A1解消 & A02持有 $e_\sigma(M)=0$ 的构造链 & 识别准确smooth type，证明全部Einstein不存在，或证明某个Einstein存在\\
Dehn-twist曲面丛 $E_D$ & 累积稿与A03第68--69步报告全局近负Einstein输入 & 输入原稿核准；非适配、非K\"ahler的一般Einstein存在性判断\\
含纯双曲块的extended graph流形 & 在Di Cerbo精确范畴内有全部Einstein不存在性 & 全Ricci缺陷趋零的闭合连接器与全局度量\\\bottomrule
\end{longtable}
曲面丛的 $\mathsf P_-$ 与图流形的 $\mathsf N$ 属于不同对象，不能相加。一个图流形局部free-$T^2$区域也不意味着整个闭流形具有自由环面作用。第三主线应保留这些差异，而不是用同一个“长颈”名称将不同拓扑合并。

\section{图流形基线与上一版编辑者的可选饱和约束}
\subsection{先固定Di Cerbo范畴，不能只用名称}
采用Di Cerbo文中定义的闭、定向extended graph四流形，检查真实双曲块、非纯块、边界仿射粘合及周边子群的条件\cite{DiCerbo}。不能把任意有图式分解的流形都包含进去，也不能将无非平凡拼接的闭双曲流形当成无Einstein例。

对于含至少一个纯实双曲四维块的相应对象，记各块截面曲率 $-1$ 的体积和为
\[
V_0=\sum_j\Vol_{-1}(P_j)>0.
\]
该文Lemmas 8--9给出 $\tau(M)=0$、$\chi(M)=3V_0/(4\pi^2)$，以及 $\Ric\ge-3g$ 规范下的最小Ricci体积为 $V_0$。它们与全局不存在性证明共同构成可靠基线。没有纯块的 $\chi=\tau=0$ 情形本来就受紧夹逼的Chern--Weil障碍限制，不应作为本路线的主要候选。

\subsection{近负Einstein序列必须同时饱和体积与Weyl预算}
以下为\status{上一版编辑者补充}的文献推论，并非原归档已经核准的新增成果。本次仅保留以供比较，不以它替代原Step100相对CGB任务，也不把它强制指定为优先路线。改用原题的 $\Ric\approx-g$ 归一化，令
\[
V_*=9V_0=12\pi^2\chi(M).
\]
若 $\delta=\delta_-(g)\le1/3$，则
\begin{equation}\label{eq:graphvol}
\boxed{\frac{V_*}{(1+\delta)^2}\le\Vol_gM
\le\frac{V_*}{(1-\delta)^2}.}
\end{equation}
左界来自最小Ricci体积：$\Ric_g\ge-(1+\delta)g$，缩放 $h=(1+\delta)g$ 后 $\Ric_h\ge-h$，其四维体积为 $(1+\delta)^2\Vol_gM$。右界来自CGB与下述点态估计。

对四个Ricci特征值处于 $[-1-\delta,-1+\delta]$，记 $\lambda_i$，则
\[
\frac{s^2}{24}-\frac{|E|^2}{2}
=\frac{(\sum_i\lambda_i)^2}{6}-\frac12\sum_i\lambda_i^2
\ge\frac23(1-\delta)^2,\qquad\delta\le1/3.
\]
可逐变量利用该二次式的凹性把最小化归到区间端点；按高端点数目 $j=0,\dots,4$ 检查五个值即可。结合式\eqref{eq:CGB}和 $8\pi^2\chi=\frac23V_*$，得到式\eqref{eq:graphvol}，并进一步得到
\begin{equation}\label{eq:graphW}
\boxed{0\le\int_M|W|^2\dd\le8\pi^2\chi(M)
\left[1-\left(\frac{1-\delta}{1+\delta}\right)^2\right].}
\end{equation}
所以固定图流形若存在 $\delta_i\to0$，必须有 $\Vol_{g_i}M\to V_*$、$\int|W|^2\to0$；变化拓扑时则是相对 $\chi(M_i)$ 的Weyl能量趋零。另有 $\int|E|^2\le4\delta_i^2\Vol M_i$。

这个约束比“有可能通过塌缩逃走”具体得多。但它不是矛盾：局部极限可以是真双曲几何，而真双曲Einstein局部片本来存在。不能把A01的\emph{禁止}校准谱论证照搬到这里。图流形的难点是，非双曲的粘合拓扑能否被藏进很小的Weyl能量与受控Ricci缺陷中。

\begin{target}[图流形的二择一结论]
在一个完整核准的上述图流形或明确序列上，构造 $\delta_-(g_i)\to0$；或者建立阻止式\eqref{eq:graphvol}--\eqref{eq:graphW}同时逼近的统一几何代价。固定流形正间隙、固定拓扑族正间隙和仅依维数的正间隙必须分别陈述。
\end{target}
若只证明模型对称类内的正间隙，不能说所有图流形度量已被排除；若真的构造成功，因为外部全部Einstein不存在性已核准，将得到负分支反例，但完整 $Q_4$ 的否定仍需正分支。

\section{第100步留下的真正全局预算：相对Chern--Gauss--Bonnet}
在本节临时使用双曲截面曲率 $-1$、$\Ric=-3g$ 的规范，记物理缺陷 $\mathsf E=\Ric+3g$、$e=\tr\mathsf E$。转回原题取 $h=3g$，则 $\delta_-(h)=(1/3)\norm{g^{-1}\mathsf E}_\infty$，不要混用数值阈值。

设 $N$ 是替换双曲cusp尾部的全部区域，删除尾部总体积为 $V_{\rm rem}$，$V_N=\Vol N$。假设补集保持原精确双曲度量，且拓扑Euler预算对应原块总和。全局CGB减去补集后有
\begin{equation}\label{eq:relative}
\boxed{
2\int_N|W|^2+12(V_N-V_{\rm rem})
=\int_N\left(|\mathsf E^\circ|^2+2e-\frac{e^2}{12}\right).}
\end{equation}
这里必须是匹配后的光滑闭度量；若只对尚有边界的试验域使用，则应保留对应边界转移项。

若端同态误差特征值 $x_i\in[-\epsilon,\epsilon]$，二次部分为
\[
\sum_i x_i^2-\frac13\Big(\sum_i x_i\Big)^2,
\]
其绝对值不超过 $4\epsilon^2$，而 $|2e|\le8\epsilon$。所以
\begin{equation}\label{eq:relativebound}
\left|2\int_N|W|^2+12(V_N-V_{\rm rem})\right|
\le(8\epsilon+4\epsilon^2)V_N.
\end{equation}
近零缺陷必满足
\[
V_{\rm rem}-V_N=\frac16\int_N|W|^2+o(V_N).
\]
\source{A03第100步修正后的相对CGB关系；本稿重新检查规范与代数。它取代第99步错误地将局部free-$T^2$连接器当成整个闭 $\chi=0$ 流形的论证。}

它是一条\emph{相容条件}，不是自动不可能性。$V_N$可以小于 $V_{\rm rem}$；体积差可以支付Weyl能量。满足该条件也没有解决水平、形状、混合和荷方程。第99步将目标改成固定缺陷约1.098的主张已撤回：任一有限正缺陷都不能反驳允许任意窄区间的原题，该数值也没有被证明几何最优。

\section{连接器方程：只保留完整的几何问题}
\subsection{Gaussian表示与全部Ricci分量}
令 $g=dr^2+q(r)$，$K=\frac12\partial_rq$、$S=q^{-1}K$、$H=\tr S$。在 $\Ric\approx-3g$ 规范下，残差的全部分量为
\begin{align}
\mathsf E_{rr}&=3-\tr(q^{-1}K_r)+\tr(S^2),\\
\mathsf E_{ri}&=q^{jk}\nabla_jK_{ki}-\partial_iH,\\
\mathsf E_{ij}&=\Ric(q)_{ij}+3q_{ij}-(K_r)_{ij}-HK_{ij}
+2(Kq^{-1}K)_{ij}.\label{eq:fullRic}
\end{align}
等价的法向方程是
\[
H'=3-|S|^2-\mathsf E_{rr}.
\]
混合Codazzi块不能省略。计算范数需通过 $g=LL^{\mathsf T}$ 的Cholesky分解，使用 $L^{-1}\mathsf E L^{-\mathsf T}$ 的谱，而不是坐标分量最大值。

中心Cauchy数据 $(q_0,K_0)$ 应满足标量与动量约束。精确Einstein规范的标量约束为
\[
R_{q_0}=-6+H_0^2-|S_0|^2.
\]
全测地中心 $K_0=0$ 要求 $R_{q_0}=-6$；最小中心 $H_0=0$ 并不要求 $K_0=0$。不能把charged中心擅自替换成全测地中心。

\subsection{正定矩阵表示不意味着方程已解}
第98步保留了全部轨道矩阵和混合方向。令 $B=q_0^{1/2}$，$A=B^{-1}K_0B^{-1}$，$\widehat R=B^{-1}\Ric(q_0)B^{-1}$，写
\[
q=B\exp(X)B.
\]
在 $H_0=0$ 的中心，两侧满足
\[
X(0)=0,\quad X_r(0)=\pm2A,\quad X_{rr}(0)=2(\widehat R+3I).
\]
因 $\frac{d^2}{dr^2}e^X|_{X=0}=X_{rr}+X_r^2$，无须假设矩阵交换，就恢复正确二阶Einstein中心数据。指数表示保证正定，却不控制全Ricci残差，也不证明可延拓到远端。

数值参数化必须\emph{解析地}固定端点一、二阶导数。若只是评估时覆盖端点导数数组，优化结果并不对应实际光滑度量。第98步粗网格中心Codazzi误差曾被误当最终缺陷；其随Fourier分辨率下降说明这是离散误差，不是几何下界。一次拟合下降、未保存系数、未去混叠的非线性乘积，不构成存在证据。

\subsection{平切片障碍与Koiso局部可解性分别说明什么}
若所有切片平坦，$R_q=0$，标量约束迫使
\[
H^2-|S|^2=6+O(\epsilon),\qquad
H^2\ge9-6\epsilon.
\]
足够小误差时无法穿过 $H=0$。这只排除平切片模型，不排除曲切片、twist或非对角补偿。

另一方面，取非恒定解析周期 $f$，
\[
h_0=d\theta^2+e^{2f(\theta)}dx^2+e^{-2f(\theta)}dy^2,
\qquad R(h_0)=-2(f')^2.
\]
负Yamabe类的归一化代表 $h=\psi^4h_0$ 可满足 $R(h)=-6$，对应
\[
-8\psi''-2(f')^2\psi=-6\psi^5.
\]
它提供Koiso式解析局部Einstein Cauchy数据，说明不能用一个失败的有限参数中心模型排除所有局部中心。它不提供双端cusp匹配、共同长时间延拓或全局闭度量。

\section{长颈、周边格点与荷：不能只优化一个变量}
\subsection{半长与误差的共同极限}
在最小中心 $H(0)=0$ 与扩张端 $H(R)=3$ 的受限Gaussian模型中，Raychaudhuri比较给出
\[
1\le k\tanh(kR),\qquad k=\sqrt{1+\epsilon/3}.
\]
若 $\epsilon\to0$，必须 $R\to\infty$；展开相应必要代价为 $\epsilon\gtrsim12e^{-2R}$。这是该端点条件下的法向必要条件，不是全连接器的充分误差率。增加 $R$ 后切向与混合残差可能更坏，不能只展示法向块衰减。

精确外部基准 $q=e^{2r}q_0$ 的残差为切向 $\Ric(q_0)$，物理大小随 $e^{-2r}$ 衰减；最大值仍可能在腰部。$q=\cosh^2(r)q_0$ 也不能保留一般charged中心的Weingarten数据。因此外部redshift不能自动修复内过渡。

\subsection{周边覆盖的正确作用}
轨迹第69、100步使用周边格点 $(N,1,1)$，经cusp深度约 $\log N$ 后得到长度尺度 $(1,N^{-1},N^{-1})$，与rank-two商塌缩中心匹配。商操作在局部保留度量与Ricci端同态，却把体积与缺陷积分同时缩小；它不直接减小点态缺陷。

第100步提出一个24-cell双曲块的具体候选及一cusp循环覆盖。该census标签、周边同调像、覆盖连通性、cusp类型和Di Cerbo假设都要回到原构造或可复核的三角剖分核准；本稿不将轨迹中的census数值当成外部已认证输入。原始周边字符非零并不自动意味着任意度数覆盖都只有一个cusp，应检查相应指数与互素条件。

对固定无量纲连接器轮廓，体积、Weyl能量与所有缺陷积分可共同按 $N^{-2}$ 缩放，式\eqref{eq:relative}除以体积后仍不变。真正近零缺陷时 $R$也要增长，匹配深度变为
\begin{equation}\label{eq:RN}
T=\log N-R+O(1),\qquad N\gg e^R.
\end{equation}
所以需要同时控制 $R,N$ 与依 $R$ 变化的轮廓体积和谱常数。把 $N\to\infty$ 当成一切误差自动消失的极限是错误的。

\subsection{荷与Bochner的抵消必须保留}
第95步的有效工具可抽象为一条带边界项的代数恒等式。对对称矩阵 $A$、向量荷 $q$、密度 $d\mu_\sigma$，设 $Q=\frac12q^{\mathsf T}Aq$、$m_q(X)=q^{\mathsf T}A\,Xq$，则
\begin{equation}\label{eq:charge}
\int m_q(X)d\mu_\sigma=
\int_{\partial\Omega}Q\langle X,\nu\rangle d\mu_{\sigma,\partial}
-\frac12\int q^{\mathsf T}\{X(A)+(\operatorname{div}_\sigma X)A\}q\,d\mu_\sigma.
\end{equation}
它由 $XQ=\frac12q^{\mathsf T}X(A)q+m_q(X)$ 分部积分得到。应用于连接器时应恢复原稿 $A=e^{-3\sigma}M^{-1}$ 的变量字典和边界通量。

原稿曾试图把drift曲率的 $+Q|dU|^2$ 当成有利项，直接沿用无twist上界；实际张力项产生 $-Q|dU|^2$ 抵消，剩余含面积梯度与形状梯度的相关项，没有固定符号。这个抵消保留了twist参与过渡的可能性。不能再使用已撤回的scalar-charge ceiling，也不能因为“有荷可能补偿”就认定全局解存在。

\section{曲面丛：已有输入之后，怎样研究一般Einstein存在性}
累积稿第6章及A03第68步记录了一类Dehn-twist曲面丛 $E_D$ 的最大范数近负Einstein输入构造，使用慢变度量、显式Ricci块、全局修正与长颈估计。当前归档没有完整给出所有种数、单值化映射与系数；续接须取回原定义，不以一个同名抽象曲面丛替代。首先确认是固定smooth type上的 $e_-(E_D)=0$，还是随 $D$ 变化的序列，这影响原题量词。

适配某个纤维化、水平可积或基本平均曲率的Einstein刚性结果，可以限制\emph{适配度量}，却不能限制任意Einstein度量。曲面丛拓扑本身不要求未知Einstein度量尊重原纤维结构。

对基面NPC目标的能量极小调和映射，不自动是调和态射或水平共形。将两个双曲曲面的乘积投向另一个标记的双曲曲面，便可产生Einstein源上的调和映射而缺少所需共形性。优化目标Teichm\"uller结构给出的通常是积分TT平稳条件，不是逐点Hopf微分消失。欲获得全局障碍，应从具体无限Dehn-twist单值化、谱退化或共形缺陷找到新的强制量。

可以考察以下有数学裁决力的方向：是否某个任意Einstein度量都诱导特定可控的调和映射结构；是否可由单值化推出一个与Einstein方程不相容的积分恒等式；或是否存在不同复结构/不同构造实际上给该smooth type一个Einstein度量。每一方向都需要完整假设，而不是“曲面丛没有复结构所以没有Einstein”。

\section{全局变分方法：保留下确界，撤掉非达到性的过强推断}
四维有
\begin{equation}\label{eq:F}
\mathcal F(g)=\int_M(|W|^2+s^2/24)\dd
=8\pi^2\chi(M)+\frac12\int_M|E|^2\dd.
\end{equation}
如果某个\emph{固定}流形的输入构造还给出体积控制，使 $\int|E|^2\to0$，则其下确界为 $8\pi^2\chi$。若光滑度量达到这个下确界，则 $E=0$，由收缩Bianchi在 $n>2$ 得标量常数，从而是Einstein度量，符号不限。

但展示一列退化极小化序列，只证明一条极小化路径不紧，\emph{不证明不存在另一光滑极小值点}。近Einstein的 $L^p$ 残差也不是四阶变分Euler--Lagrange残差。Ekeland、Palais--Smale或“泛函非负”不能未经估计填平此差距。每次调用必须说明固定范畴、紧性、规范、允许退化和最终smooth type的恢复。

一个真正有影响的结果可以是：在这些已构造输入的流形上证明全局达到性，产生Einstein；或者证明所有可能达到情形都被一个严格障碍排除，得到全部Einstein不存在性。任一方向都要覆盖\emph{全部候选极小值}，而不是只覆盖选定orbifold附近的一支。

\section{不能遗漏的正分支与高维问题}
本会话不重复A02的有限维计算，但应维护全局覆盖范围。$\Ric\approx+g$ 给Myers直径界、Bishop体积上界和有限基本群，不给统一体积下界或基本群阶上界。精确Einstein商可以塌缩；正确目标不是排除所有塌缩，而是理解其是否仍与Einstein存在性兼容。

四维在固定非塌缩条件下可用有界Ricci有限型与正则片理论\cite{CN}。有限多个微分同胚型本身不构造Einstein度量，也不使极限光滑；曲率集中与smooth-type恢复仍需研究。对两两不同的高拓扑复杂度输入，体积塌缩可能被迫出现。此处的四维曲率能量约束与A01的SW近等号密度不同，不能强行套同一禁止谱。

若在一个非塌缩覆盖上产生了Einstein度量，仍要证明它对原甲板群等变，才能下降；覆盖存在Einstein并不自动意味着商的指定光滑结构存在所需度量。

对于 $n\ge5$，四维的 $\Lambda^2_\pm$、SW/Bauer--Furuta和CGB二次能量机制没有直接同构替代。高维路线可以寻找拓扑驱动的临界 $L^{n/2}$ 预算、特定特殊结构下的禁止曲率模式、沿余维四奇异层的ALE束障碍或新的全部Einstein不存在性。但若加入特殊holonomy、对称性、体积或完整曲率假设，只得到该范畴定理，不是原题。

四维输入与同符号Einstein因子相乘，端同态误差是两因子误差的最大值，因而可传播某些\emph{输入}；它不传播全部Einstein不存在性，因为产品上的任意Einstein度量不必分裂。正Einstein球因子维数至少2，这种乘积也跳过五维。项目尚无可直接调用的高维全部Einstein障碍，不等于数学上不存在，也不构成停止探索的理由。


\section{其他高价值路线的保留方式，而非重新从零铺开}
双曲连通和与BCG/Sambusetti熵路线在轨迹中已有目标依赖的障碍信息；不足处是没有与之匹配的任意紧Ricci夹逼构造，且常数可能依赖目标。其正确延伸是研究近等号熵映射能否保留周边同伦信息、能否给出固定拓扑损失的定量代价，而不是再次把经典熵不等式本身列为新突破。使用前恢复A01第91步的实际命题与源文献，不能把某个目标的界推广为维数统一界。

精确Poincar\'e--Einstein或AH/cusp模型显示切向与法向非线性补偿并非不可能，这是反对简单局部无解结论的真实依据。但一个自对偶模型未必提供可反射的全测地 $T^3$ 中心：若解析Einstein邻域中有反射，该反射逆转取向并交换 $W^+,W^-$，自对偶与反射相容会迫使二者同时消失；中心因而不能是所设想的紧致负常曲率三环面。这只限制该反射应用，不否定其他非对称中心。

已有负Einstein Dehn filling思路可以帮助研究边界与逆算子，但填充改变周边拓扑；得到另一个流形上的Einstein度量，不能成为Di Cerbo图流形上的存在反证。有限维的free-$T^2$、对角、orthogonally transitive或反射对称类也不穷尽所有图流形度量。保留这些路线的价值在于给主攻问题提供精确对照，而不是用局部范畴的失败宣布原题已经偏向真或假。

\section{本会话应优先裁决的三个具体问题}
\subsection{问题G：真实双端连接器的连续极限是否存在}
固定完整Di Cerbo图流形候选、周边粘合和中心约束，构造保持两侧端点全部jets的完整矩阵场。数值上要求系数保存、精细网格独立验证、非线性去混叠、正定性下界、所有Ricci块和尺度依赖误差；再寻求解析或区间后验估计。若线性化逆范数为 $B$、二次非线性常数为 $C$、残差为 $\varepsilon$，校正可行性至少要控制相应 $B^2C\varepsilon$，而不是只让离散残差变小；余核与自由参数需另外解决。

若始终失败，寻找在分辨率和端点参数变化下仍存在的连续伴随障碍，而不是把一次优化停滞当成不存在性。有效障碍须对真实方程与允许自由度配对，并能排除所有足够小缺陷，而非只排除某个五次插值。

\subsection{问题S：曲面丛或固定解消的全局Einstein判断}
优先挑一个输入已经最清楚的具体smooth type，恢复全部拓扑数据和构造证明。对正、零、负Einstein常数逐项列出可用证书，剩余的某个符号不能用未经证明的Ricci-flat分类消掉。尽可能把问题转成一个具实际约束的调和映射、能量达到性或全局模空间命题。

如果证明了Einstein存在，应明确是新存在性结果而非反例失败的废弃产物。若只证明附近无解，保留为局部定理，并说明何种额外全局覆盖才能得到全部不存在性。

\subsection{问题R：图流形的Weyl小量是否有全局拓扑代价}
式\eqref{eq:graphW}使假想序列具有归一化Weyl小量。研究这种\emph{全局}近双曲状态能否强迫相容的几何/拓扑结构，或允许非双曲周边 $\Z^3$ 在极小能量区域存活。不能把 $L^2$ 小量直接升为全域点态Weyl小量；A02的尖峰是必要测试。

若能构造一个新的定量拓扑稳定性定理，可能直接排除整类负分支反例；若找到能保留周边群的严格构造，可能完成问题G。它与A01有共同的“小缺陷到局部结构”技术，但这里的局部双曲模式\emph{不是被禁止的模式}，最后必须有全局步骤。

\section{Prior-art边界与论文可能的核心}
Di Cerbo的全部Einstein不存在性和最小Ricci体积不是本项目成果\cite{DiCerbo}。由它们结合CGB得到式\eqref{eq:graphvol}--\eqref{eq:graphW}是本文件重算的直接推论，可用作严格筛选条件，不以一个代数推论单独宣称重大突破。Koiso局部解析存在、Yamabe归一化和负曲率调和映射理论也是背景；真正新增要落在全局匹配、近零缺陷的共同参数控制或无额外适配假设的Einstein判断上。

当前不承诺这条主线已经有一篇完成的主论文。可独立成文的强结果应当是：一个全局负分支反例；一个无对称性限制的连接器/图流形夹逼间隙；一个此前未知smooth type的Einstein存在性；或一个显著更一般、可用于多个候选的全局稳定性/非达到性定理。仅若某受限模型定理本身具有独立新意与用途，才考虑将其单独写出。

\section{接续入口与不能丢失的原始资产}
\asset{第100步相对CGB与全局预算}{Sessions/S0910A03/workspace/notes/step100_relative_cgb_graph_budget.md}
\asset{全矩阵匹配压力测试}{Sessions/S0910A03/workspace/notes/step98_full_matrix_matching_pressure_test.md}
\asset{荷与Bochner抵消}{Sessions/S0910A03/workspace/notes/step95_KX_and_bochner_cancellation.md}
\asset{周边格点与各向异性匹配}{Sessions/S0910A03/workspace/notes/step69_anisotropic_cusp_compensation.md}
\asset{局部Koiso中心数据}{Sessions/S0910A03/workspace/scripts/step69_collapsed_koiso_center.py}
\asset{完整矩阵残差程序}{Sessions/S0910A03/workspace/scripts/step98_full_matrix_joint_transition.py}
\asset{独立分辨率检查}{Sessions/S0910A03/workspace/scripts/step98_center_residual_resolution.py}
\asset{曲面丛、全局TT与单值化章节}{Skills/qiushi-yau26/references/research/20260912.tex}
\asset{双曲熵障碍的已有路线}{Sessions/S0910A01/workspace/notes/step91_sambusetti_bcg_quantitative_route.md}
\asset{同伴固定A1输入}{Sessions/S0910A02/workspace/notes/step83_near_pinched_construction_and_frontier_reduction.md}
\asset{同伴精确拓扑间隙}{Sessions/S0910A03/workspace/briefing/Yau26_uniform_gap.tex}

第99步笔记与JSON应带第100步修正，不作为当前反例目标。第98步的端点数组覆盖、粗网格最大值与未保存系数不得用于严格构造。第95步charged ceiling已撤回，但式\eqref{eq:charge}及真实抵消应保留。第68--69步曲面丛调和态射缺口、非适配Einstein自由度与周边匹配问题不得被后续摘要省略。

\subsection{从Step101开始的合理顺序}
先读取当前工作区与上述最终笔记，冻结一个真正满足外部不存在性假设的图流形候选及其完整周边数据；同时核准曲面丛输入的准确版本。用相对CGB、全矩阵残差与 $(R,N)$ 联动重算候选的全局预算，然后在问题G、S、R中选择一个有最强证据的目标深入。

A03不接管A01已经审读的精确间隙主证明，也不复制A02十维Weyl核计算。需要调用时传递完整条件与版本。若发现新的更有效全局机制，允许转向；但必须说明它改变了 $\mathsf P_\sigma$ 或 $\mathsf N$ 的哪一部分，不能仅增添一组相关术语。

这一安排是可调整的编辑建议，不是已穷尽300步价值的证明；原始成果与分支由覆盖表和原件继续保留。后续必须仍围绕：同一光滑对象、同一归一化、全部方程、全部Einstein符号和明确常数依赖。

\section{覆盖补全：遗漏的存在性分支、全局筛选与最终接续}\label{sec:restore-a03}
上一版把本稿过度集中在图流形和全局反例，而未充分保留已有的\emph{Einstein存在性}工作。本节恢复这些内容，来源为三份实际上传meta-trace与归档明示的条目；本轮不对未附的证明和代码再作认证。

\subsection{Q-Gorenstein平滑：一个不能被删去的存在性研究簇}
A01第49--53步研究的链条为
\[
\begin{gathered}
\text{全局Q-Gorenstein平滑}\longrightarrow\text{canonical ample保持}\\
\longrightarrow\text{Aubin--Yau负KE}\\
\longrightarrow\text{simultaneous resolution/Ehresmann}\\
\longrightarrow\text{指定smooth resolution上Einstein}.
\end{gathered}
\]
它可能将原反例候选转化为存在性例子。只证明局部节点有平滑方向，不等于全局平滑；中央复结构非ample，也不等于底层光滑流形无Einstein。

第51步\emph{报告}对BCGP的 $G=\Z_4^2$、$K^2=2$、六个 $A_1$ 节点行完成了存在性证明：通过
\[
T_X^0\simeq(\pi_*T_Y)^G,\qquad
H^2(X,T_X^0)=H^2(Y,T_Y)^G,
\]
及等变Serre对偶、K\"unneth与特征标计算使两项不变空间消失，再用全局平滑、ample与微分同胚桥。其原稿为\path{Sessions/S0910A01/workspace/notes/step51_reflexive_bridge_einstein_existence_proof.md}。17/17检查是源记录的附带证据，不代替整条几何证明；本次未取得原稿，故准确状态是\emph{轨迹报告了闭合个案，需恢复原证明与查新}，不是“从未有结果”，也不是本次已重新验证世界首创。

第50步的经典Campedelli构造使用七条直线
\[
x,y,z,x-y,y-z,x-z,x-y+z=0
\]
并检查三重点标签；不能仅因基本群或 $K^2$ 相同就把它与BCGP行认成同一平滑分支。第52步作了 $K^2=4,6$ 多行的holomorphic Lefschetz obstruction扫描；第53步在 $\mathbb Q(\zeta_{60})$ 精确检查BCGP \S6.9的 $d_{21}=d_{12}=1$、四个节点轨道。即使 $\dim\ker\beta\ge2$，仍不推出存在所有节点坐标都非零的平滑方向。这些具体结果和缺口应共同保留。

此研究簇由本稿组织承接，但仍属于原A01已开展的工作；这不是把功劳或实际工作区转移给A03，也不要求A01丢弃原分支。

\subsection{全局工具为何与近Einstein输入错位}
A02第86--92步的来源审计依次处理PSC与monopole不相容、反向 $A_1$ 填充中的 $+2$ 球对具体BF资料的限制，以及ALE Euler/signature预算与BCG熵界的关系。BF论证必须保留closed-to-relative接口、非零性、取向和嵌入球条件；不能从“不适用这个BF证书”断言没有任何全局障碍。

A01第91--92步的 $N\#N$ 双曲／Sambusetti--BCG路线不应只留一句“目标依赖”。它具有明确的无Einstein或度映射障碍对象，缺的是对应的任意紧点态输入及统一性；其阈值必须保留目标、直径、次数等真实依赖。原入口为\path{step91_sambusetti_bcg_quantitative_route.md}及第92步修正版。A02第90步熵推论曾在第92步修正缩放方向：$\widehat g=(1+\epsilon/3)g$，四维体积乘 $(1+\epsilon/3)^2$，不能使用倒数版本。

相应资料也保留了ALE添加对 $\chi-\frac32|\tau|$ 的影响分析；它说明某些经典障碍不能直接与给定近似构造配对，\emph{不等于全局Einstein存在性已经证明}。粗alignment、未知Ricci-flat分类、某个复结构的K\"ahler限制，都不能代替缺失的全符号排除。

\subsection{图流形末步的具体对象、数值证据和原样待办}
A03第100步并非只有相对CGB公式。其引用的候选块为\texttt{24cdc1.1}，源记录列出：可定向单 $T^3$ cusp、$\chi=2$、截面曲率 $-1$ 体积 $8\pi^2/3$、最大cusp体积16、$b_1=1,b_3=0$。\emph{这是轨迹引用的外部数据，本次未重新核准该三角剖分或原文}；但为了不丢失具体研究对象，必须原范围保留，而不是只写“某个24-cell”。

原第100步要求：核准周边 $H_1$ 像及primitive character、互素覆盖度子序列、一cusp条件、两块取向相容的仿射粘合和Di Cerbo假设；同时计算 $R,N$ 联动时的真实Weyl、体积和所有Ricci块。已有JSON的connector/Weyl密度是\emph{占位参数}，$N^{-2}$缩放核对不是一个已测量的实际连接器。旧Step99的JSON未重生成，不能继承其连接器解释。

A03第69步还记录Koiso数值演化的分辨率梯队：6/8/10/12模在早期约 $t\lesssim0.4$ 相互趋近，并出现 $H$ 先正后负；更晚阶段刚性与分辨率依赖增强，不能当作几何爆破证明。该初始族的 $S'(0)$ 与互补截面曲率关系只限制相应反射中心。它不是一般双端Einstein连接器的不存在性定理。

第94--96步的带荷中心、水平Noether配对、cusp indicial根和全局残差map也应保存。第96步报告全系统主符号rank六，及特定规范中的scalar根 $-1,0,3$、orbit根 $0,3$；这些是\emph{原模型的局部分析资料}，不自动给全局Fredholm逆或双端解。第93步曾因在最大值处 $LQ<0$ 误判矛盾，第94步已撤回；第95--96步又修正了 $K_X$ 为恒等配对而非新障碍的解释。

原入口另包括\path{step94_charge_sign_correction_and_global_flux.md}、\path{step96_charged_koiso_local_evolution_and_joint_matching.md}、\path{step100_relative_cgb_scaling.py}和\path{step100_relative_cgb_scaling.json}。这些路径与原最后待办已完整保存在300步索引，不以本版路线建议替换。

\subsection{曲面丛：存在的输入、失败的推断与尚缺的原始定义}
$E_D$ 的全部种数、单值化、度量系数和误差常数在当前ZIP中未给全；不能自行补全这个对象。归档保存的核心内容是最大范数近負Einstein输入、长collar的逆估计问题，以及一般Einstein度量不必适配给定纤维化。

具体地，$L^2$谱控制不自动提供长度一致的 $L^\infty$逆；近零模和Green核可能随颈长放大。A03第68--69步已用乘积双曲例及目标Teichm\"uller优化检查，否定“调和、submersion、唯一能量极小就推出水平共形／调和态射”的捷径。须保留这些明确失败机制，不能只留下“曲面丛还缺全局障碍”的一句话。

\subsection{一般正分支与高维必须有明确接续位置}
本版将一般四维正分支的非塌缩集中、塌缩、单连通条件性体积归约恢复在A02；它不是A03可无限期代管的尾注。A03负责各构造的same-smooth-type判断与全符号存在性，不承担替代所有维数理论的任务。

归档09保留的高维关键观察为
\[
A_U\|\Rm\|_{L^{n/2}(U)}\lesssim r^2
\left(\fint_U|\Rm|^{n/2}\right)^{2/n},
\]
这是其相应Sobolev/能量约定下的尺度结构，\emph{不是任意塌缩域免费获得Sobolev不等式}。待补的是拓扑驱动的全局预算和禁止局部曲率模式，而非仅把四维指数换成 $n/2$。

$n=5$、余维四ALE束、群表示不变量、全维正分支与新的全部Einstein不存在性，均应保留为证据驱动的开放分叉；不因为当前两条成果成熟就宣布这些路线没有价值。乘积只传播已构造输入，不传播任意Einstein非存在性。原题完整解决的最高目标不因分成三份续接稿而缩小。

\section{结果交付与论文标准}
研究结论以具体命题和证明为单位，而不是以Agent数量、步数或平台认可次数为单位。同一公式在多个会话复算，可以提高证据质量，不产生多项独立新发现。

可完成的结果应及时冻结成可独立阅读的数学稿：准确的假设与量词，新增结论相对最接近文献的差别，所有分析接口及常数依赖，完整证明，必要时附有限计算及独立复算说明。假设增多或结论缩小必须在主定理中显式反映。无需等到整个Yau26解决才允许发表严谨的中间结果，也不把一个证明链拆成若干缺乏独立贡献的论文。

Annals、Inventiones、JAMS、Acta可以作为高质量写作与结果强度的目标，但任何文件都不预设已达到其录用标准。完整性、首创性、重要性和选刊是四项判断。用经典工具证明新的重要定理可以完全原创；出现一个新名称、搜不到同名词或构造了“框架”，并不能确认首创性。按\emph{精确命题及假设}进行查新，并记录最近定理、包含关系及真正新增的证明步骤。

无论本轮选择论文收束还是进一步探索，结束时只需留下四类必要材料：更新后的主定理与证明；当前最小未闭合命题；确切被否定的模型及原因；可供同伴直接使用的来源、版本和计算。除非出现新的反证或新缺口，不再次从零争论已审读的同一结果。反过来，发现具体错误就修正所有依赖它的结论，而不是只新增一份不被后续读取的纠错说明。

\begin{thebibliography}{99}
\bibitem{Yau1982} S.-T. Yau, \emph{Problem Section}, in \emph{Seminar on Differential Geometry}, Annals of Mathematics Studies 102, Princeton University Press, 1982, pp. 669--706. 历史题页仍需在工作区逐字核对；本稿采用任务书中的区间命题。

\bibitem{DiCerbo} L. F. Di Cerbo, \emph{Extended graph 4-manifolds, and Einstein metrics}, arXiv:2106.13279. \url{https://arxiv.org/abs/2106.13279}. 本轮核对其 extended graph 定义、Lemmas 8--9 和 Einstein 不存在性证明；不推广到所有称为“图流形”的空间。

\bibitem{CN} J. Cheeger and A. Naber, \emph{Regularity of Einstein manifolds and the codimension 4 conjecture}, arXiv:1406.6534. \url{https://arxiv.org/abs/1406.6534}. 引用须区分有界 Ricci 的非塌缩正则片与精确 Einstein 的更强正则性。

\bibitem{NZ} A. Naber and R. Zhang, \emph{Topology and $\epsilon$-regularity theorems on collapsed manifolds with Ricci curvature bounds}, Geom. Topol. 20 (2016), 2575--2664. \url{https://doi.org/10.2140/gt.2016.20.2575}. 最大纤维化幂零秩、模型几何等输入不能省去；本稿不宣称这些输入已由夹逼推出。

\bibitem{JKX} W. Jiang, L. Kong and S. Xu, \emph{Convergence of Ricci-limit spaces under bounded Ricci curvature and local covering geometry I}, arXiv:2205.00373. \url{https://arxiv.org/abs/2205.00373}. local covering geometry 是额外结构；具体使用前核对版本与定理。

\bibitem{Anderson1989} M. T. Anderson, \emph{Ricci curvature bounds and Einstein metrics on compact manifolds}, J. Amer. Math. Soc. 2 (1989), 455--490. \url{https://doi.org/10.1090/S0894-0347-1989-0999661-1}. 本轮确认文献身份，未取得所需 Corollary B(4) 的完整原文；不得据二手概述填写其假设或用它排除输入曲率尖峰。

\bibitem{GJ2017} H. Ge and W. Jiang, \emph{$\epsilon$-regularity for shrinking Ricci solitons and Ricci flows}, arXiv:1707.05511. \url{https://arxiv.org/abs/1707.05511}. 本轮核对 Theorem 1.1、Theorem 3.4；后者为有界 Ricci 的尺度选择，不是一般有界 Ricci 的点态曲率结论。

\bibitem{OzuchII} T. Ozuch, \emph{Noncollapsed degeneration of Einstein 4-manifolds, II}, arXiv:1909.12960v2; Geom. Topol. 26 (2022). \url{https://arxiv.org/abs/1909.12960}. 本轮重点核对 \S7.3、Corollary 15、Remarks 50--51 及明确的 $L^\infty$ 问题；编号以该 arXiv 版本为准。

\bibitem{LeBrun2001} C. LeBrun, \emph{Ricci curvature, minimal volumes, and Seiberg--Witten theory}, Invent. Math. 145 (2001), 279--316. \url{https://arxiv.org/abs/math/0003068}. 本轮核对 Theorem 2.3 的证明及式(11)、Proposition 3.1--3.2、Corollary 3.4；不把 $1/3$ 说成已证明最优的 Einstein 拓扑阈值。

\bibitem{Project} Qiushi Engine 项目资料：2026-09-15进展包、2026-09-16逐步归档，S0910A01/A02/A03各100步 meta-trace；源文件名与本地工作区定位见正文。原稿、脚本与摘要须分别识别，不能把记录中的成功标签视作外部定理。

\bibitem{CodexReview} 用户在本会话提供的 Codex 原稿审读报告（截至2026-09-16）：\path{Sessions/S0910A03/workspace/briefing/Yau26_uniform_gap.tex}，重点位置原稿约第277行；另核查 A02 第83步构造笔记。结论为在所审读依赖链中未发现实质缺口，且精确拓扑间隙不依赖(SC)。本引用是审读报告的来源定位，不声称原稿全文已包含于归档。

\bibitem{CoverageAudit} 本版覆盖复核资料（2026-09-16）：\path{审计/研究条目覆盖表.tsv}、\path{审计/300步原始任务索引.json}、\path{审计/原件校验.tsv}。恢复内容来自实际上传三份meta-trace及两个归档的明确条目；并非本次重新证明全部原稿。
\end{thebibliography}
\end{document}
